% from ega1-3-fr.tex
\oldpage[I]{116}
est commutatif, $\varphi$ et $\varphi'$ provenant du morphisme
$\operatorname{Spec}(K')\to\operatorname{Spec}(K)$, $\alpha$ et $\alpha'$
correspondant à $f$. Or, $\varphi$ est injectif et il en est de même de
$\alpha'$ par hypothèse~; donc $\alpha$ est nécessairement injectif.
\end{env}

\begin{proposition}[3.5.6]
\label{I.3.5.6-fr}
Soient $f:X\to Y$, $g:Y\to Z$ deux morphismes de préschémas.
\begin{enumerate}
  \item[(i)] Si $f$ et $g$ sont radiciels, il en est de même de $g\circ f$.
  \item[(ii)] Inversement, si $g\circ f$ est radiciel, il en est de même de
  $f$.
\end{enumerate}
\end{proposition}

Compte tenu de la déf. (\hyperref[I.3.5.4-fr]{3.5.4}), la proposition revient
aux assertions correspondantes pour les applications
$X(K)\to Y(K)\to Z(K)$, qui sont évidentes.

\begin{proposition}[3.5.7]
\label{I.3.5.7-fr}
\begin{enumerate}
  \item[(i)] Si les $S$-morphismes $f:X\to X'$, $g:Y\to Y'$ sont radiciels,
  il en est de même de $f\times_S g$.
  \item[(ii)] Si le $S$-morphisme $f:X\to Y$ est radiciel, il en est de même
  de $f_{(S')}:X_{(S')}\to Y_{(S')}$ pour toute extension $S'\to S$ du
  préschéma de base.
\end{enumerate}
\end{proposition}

Vu (\hyperref[I.3.5.1-fr]{3.5.1}), il suffit de démontrer (i). On a vu
(\hyperref[I.3.4.2.1-fr]{3.4.2.1}) que
\phantomsection\label{I.3.5.7.field-valued-product-fr}
\[
  (X\times_S Y)(K)=X(K)\times_{S(K)}Y(K),\qquad
  (X'\times_S Y')(K)=X'(K)\times_{S(K)}Y'(K)
\]
l'application $(X\times_S Y)(K)\to(X'\times_S Y')(K)$ correspondant à
$f\times_S g$ s'identifie alors à $(u,v)\mapsto(f\circ u,g\circ v)$ et la
proposition en résulte aussitôt.

\begin{proposition}[3.5.8]
\label{I.3.5.8-fr}
Pour qu'un morphisme $f=(\psi,\theta):X\to Y$ soit radiciel, il faut et il
suffit que $\psi$ soit injectif et que, pour tout $x\in X$, le monomorphisme
$\theta^x:k(\psi(x))\to k(x)$ fasse de $k(x)$ une extension radicielle de
$k(\psi(x))$.
\end{proposition}

Supposons $f$ radiciel et montrons d'abord que la relation
$\psi(x_1)=\psi(x_2)=y$ entraîne nécessairement $x_1=x_2$. En effet, il
existe un corps $K$, extension de $k(y)$, et des $k(y)$-monomorphismes
$k(x_1)\to K$, $k(x_2)\to K$ (Bourbaki, \emph{Alg.}, chap. V, \S~4,
prop. 2)~; les morphismes correspondants
$u_1:\operatorname{Spec}(K)\to X$, $u_2:\operatorname{Spec}(K)\to X$ sont
donc tels que $f\circ u_1=f\circ u_2$, donc $u_1=u_2$ par hypothèse, et cela
implique en particulier $x_1=x_2$. Considérons maintenant $k(x)$ comme
extension de $k(\psi(x))$ au moyen de $\theta^x$~: si $k(x)$ n'est pas
extension radicielle de $k(\psi(x))$, il existe deux
$k(\psi(x))$-monomorphismes distincts de $k(x)$ dans une extension
algébriquement close $K$ de $k(\psi(x))$ et les deux morphismes correspondants
$\operatorname{Spec}(K)\to X$ violeraient l'hypothèse. Inversement, compte
tenu de (\hyperref[I.2.4.6-fr]{2.4.6}), il est immédiat que les conditions de
l'énoncé sont suffisantes pour que $f$ soit radiciel.

\begin{corollary}[3.5.9]
\label{I.3.5.9-fr}
Si $A$ est un anneau, $S$ une partie multiplicative de $A$, le morphisme
canonique $\operatorname{Spec}(S^{-1}A)\to\operatorname{Spec}(A)$ est
radiciel.
\end{corollary}

En effet, ce morphisme est un monomorphisme
(\hyperref[I.1.6.2-fr]{1.6.2}).

\begin{corollary}[3.5.10]
\label{I.3.5.10-fr}
Soient $f:X\to Y$ un morphisme radiciel, $g:Y'\to Y$ un morphisme, et soit
$X'=X_{(Y')}=X\times_Y Y'$. Alors le morphisme radiciel $f_{(Y')}$
(\hyperref[I.3.5.7-fr]{3.5.7 (ii)}) est une bijection de l'espace sous-jacent
$X'$ sur $g^{-1}(f(X))$~; en outre, pour tout corps $K$, l'ensemble $X'(K)$
s'identifie au sous-ensemble de $Y'(K)$ image réciproque par l'application
$Y'(K)\to Y(K)$ (correspondant à $g$) du sous-ensemble $X(K)$ de $Y(K)$.
\end{corollary}

La première assertion résulte aussitôt de
(\hyperref[I.3.5.8-fr]{3.5.8}) et de
(\hyperref[I.3.4.8-fr]{3.4.8})~; la seconde, de la commutativité du diagramme
